% ID: SEL-D0004-OB-E04
% ESTADO: cotejada. Cotejo literal independiente y revisión de composición realizados.
% Fuente: Selectividad/sel_2019_matematicas_aplicadas/Suplente Septiembre/Exámenes Matemáticas Aplicadas CCSS - A y B.pdf
% SHA-256: eecad753694b3d8cb2d427efddb00d75faae9c4ec5b027056291bc95d7d2e3da
% Página: 2; opción B; ejercicio 4.
% Transcripción de primera pasada basada en la imagen de la página original.
% Se conservan los nueve datos en su orden, incluido el valor 10 000 repetido.
% Sin solución ni modificación de datos o condiciones.

La vida útil de los filtros de las máquinas de agua por ósmosis se distribuye según una ley Normal de media desconocida y desviación típica de 2\,000 horas. En una prueba realizada en 9 máquinas elegidas al azar, se obtuvieron los siguientes resultados:
\[
9\,500\quad 10\,000\quad 8\,500\quad 10\,500\quad 16\,500\quad 10\,000\quad 12\,000\quad 14\,000\quad 17\,000
\]
\begin{enumerate}[label=\alph*),leftmargin=2.4em]
  \item \textbf{(1.5 puntos)} Calcule un intervalo de confianza al 99\,\% para la vida útil media de los filtros de las máquinas.
  \item \textbf{(1 punto)} \textquestiondown{}Cuál debe ser el tamaño mínimo que debería tener una muestra, para que el error cometido en la estimación de la vida útil media de los filtros sea inferior a 500 horas, con un nivel de confianza del 95\,\%?
\end{enumerate}
